
we have certainly grasped knowing
as unchangeable, self-same and self-sustaining,
beyond all change and beyond the subjectivity-objectivity,
which is inseparable from change.
But this insight was not yet the science of knowing itself,
but rather only its premise.
The science of knowing must still actually construct
this inner, qualitatively unchangeable being,
and as soon as it does this,
it will simultaneously create change,
the second term, as well.

The true authentic meaning of this
simultaneous double construction
of the changeable and the unchangeable
will become completely clear
only when we actually and immediately
carry out the construction,
something that belongs within the sphere
of the science of knowing itself,
certainly not in preliminary reports.
Misunderstandings about this are unavoidable at the beginning.
In order to come as close as possible to
complete accuracy right from the start,
I venture on a question that has already been raised.

On introducing the schema:
A
x y z • B-T
I said that the science of knowing stood in the point.
I have been asked whether it doesn't rather go in A.
The most exact answer is that actually and strictly
it doesn't belong in either of these
but rather in the oneness of both.
By itself A is objective and therefore inwardly dead;
it should not remain so, indeed it should become actively genetic.
The point, on the other hand, is merely genetic.
Mere genesis is nothing at all; but this is not just
mere genesis but the determinate genesis that is required by
the absolute qualitative A; [it is] a point of oneness.
Now to be sure this point of oneness can be realized immediately,
oscillating and expending itself in this point;
and we, as scientists of knowing, are this realization inwardly
(I say inwardly and concealed from ourselves).
But this point can neither be expressed nor reconstructed
in its immediacy, since all expression or reconstruction is
conceiving and is intrinsically mediated.
It is expressed and reconstructed just as we have expressed it
at this moment: namely, that one begins from A and,
indicating that it cannot persist alone,
links it to the point;
or one begins with the point and,
indicating that it cannot persist alone,
links it to A, all the while, to be sure,
knowing, saying, and meaning that
neither A nor the point can exist by itself and
that all our talk could not express the implicit truth,
but instead that the implicitly unreconstructible something
which can only be pictured in an empty and objective image is
the organic oneness of both.
Thus, since reconstruction is conceiving, and
since this very conceiving explicitly abandons its own intrinsic validity,
this is precisely a case of conceiving the inconceivable as inconceivable.
Therefore—and I put this here first as a clarification—
this is a question of the organic split into
B-T and x, y, z, as I explained at the beginning of the last lecture;
because when I speak, I must always put one before the other.
But is it actually so?
No, instead it is exactly the same stroke;
and let me add this as well:
this deeper connection must indeed itself be a result and
a lower expression of the higher one just now described.
Finally, in order to say this in its full meaning
and thereby to make your insight into the science of knowing,
and into knowing as such, much clearer:
secondary knowing, or consciousness, with its whole lawful play
by means of fixed change and the manifold
(within it or outside it),
of sensible and supersensible,
and of time and space, comes to be, in principle,
through this recently noted and demonstrated division,
taken merely as a division and nothing else.
Everything we attribute to the subject, as originating from it,
derives from this.
Because it is clear without further ado that from a particular perspective,
namely the science of knowing's synthetic perspective,
the disjunction must be just as absolute as oneness;
otherwise we would be stuck in oneness
and would never get outside it to changeability.
(Let me note in passing that this is
an important characteristic of the science of knowing
and distinguishes it from Spinoza's system,
which also wants absolute oneness
but does not know how to make a bridge from it to the manifold;
and, on the other hand, if it has the manifold,
cannot get from there to oneness.)

As scientists of knowing, we never escape
the principle of division inwardly and empirically
(by means of what we do and promote);
but we certainly escape it intellectually
with regard to what is valid in itself,
in which very regard the principle of division
surrenders and negates itself.

Or, so that I make the point at which
we have arrived even clearer:
since we actually reason as we have just been doing,
where does our reasoning stand,
if we remain exclusively in consciousness?
(“If we remain exclusively in consciousness,”
I say, because we can also lose ourselves
in the intelligible realm,
and there is even, in its place,
an art of consciously losing oneself in it.)
Manifestly in our construction
by means of the principle of division,
it stands in the place not of that which is to be
valid to the extent it is constructed
but of that which is intrinsically valid;
thus it stands wholly autonomously, as has been said,
between the two principles of oneness and separation,
simultaneously annulling both and positing both.
Thus, the standpoint of the science of knowing,
which stands still in consciousness,
is by no means a synthesis post factum;
but instead a synthesis a priori,
taking neither division nor oneness as given,
but creating both at one stroke.
Once again to adduce an even higher perspective
what is the absolute oneness of the science of knowing?
Neither A nor the point,
but instead the inner organic oneness of both.
Besides this given description of
the point of oneness
is there also another?
None at all, we have seen.
Therefore, this description is
the original and absolutely authentic one.
What are its constituents?
The organic oneness of both is
a construction or a concept,
and indeed the single absolute concept,
abstracted from nothing existing,
since even its own separate existence,
and hence the existence of everything conceptual,
is denied.
Further, the construction as such is denied
by the manifestness of what exists autonomously;
thus even the inconceivable, as the inconceivable
and nothing more, is posited by this manifestness,
posited through the negation of the absolute concept,
which must be posited just so that it can be annulled.

And so:

1. the necessary unification and indivisibility
of the concept and the inconceivable is clearly seen into,
and the result may be expressed thus:
if the absolutely inconceivable is to be
manifest as solely self-sustaining,
then the concept must be annulled,
but to be annulled, it must be posited,
because the inconceivable becomes manifest
only with the negation of the concept.
Supplement: hence inconceivable = unchangeable; concept = change.
Therefore along with the foregoing it is evident that
if the unchangeable is to appear, there must be change.

2. Now to be sure inconceivability is only
the negation of the concept, an expression of its annulment.
Therefore it is something which originates
from both the concept and knowing themselves;
it is a quality transferred by means of absolute manifestness.
Noting this, and therefore abstracting from this quality,
nothing remains for oneness except absoluteness
or pure self-sufficiency in itself.

3. The following consideration makes this
particularly important and relevant:
What is pure self-sufficient knowing in itself?
The science of knowing has to answer this question,
or, as we put it more precisely,
it has to construct the presupposed inner quality of knowing.
We are undertaking this construction here:
negating the concept by means of manifestness,
and thus the self-creation of inconceivability is
this living construction of knowing's inner quality.
This inconceivability itself originates in the concept and
in pure immediate manifestness;
likewise the whole quality of the absolute,
as well as the fact that a quality can be applied to it at all,
originates in the concept.
The absolute is not intrinsically inconceivable,
since this makes no sense;
it is inconceivable only when the concept itself tries for it,
and this inconceivability is its only property.
Having recognized this inconceivability as
an alien quality introduced by knowing, I said before,
only pure self-sufficiency, or substantiality, remains in the absolute;
and it is quite true that at best this
self-sufficiency does not originate in the concept,
since it enters only with the latter's annulling.
But it is clear that this quality enters only
within immediate manifestness, within intuition,
and thus is only the representative and correlate of pure light.
This latter is its genetic principle by which, first of all,
according to our hypothesis all manifestness opens up
into genetic manifestness, since pure light
manifests itself implicitly as genesis.
Secondly, the previously presented relationship of
concept to being and vice versa is further determined as follows:
If there is to be an expression and realization of the absolute light,
then the concept must be posited,
so that it can be negated by the immediate light,
since the expression of pure light consists just in this negation.
But the result of this expression is being in itself, period.
[This result] is inconceivable precisely because
pure light is simultaneously destruction of the concept.
Thus, pure light has prevailed as the one focus and
the sole principle of both being and the concept.

4. From the preceding it follows that this inconceivable,
as the bearer of all reality in knowing,
which we grasp in its principle, is absolute only as inconceivable,
and cannot be thought in any other way.
No other additional hidden qualities are attributable to it.
Just as little can any qual-
ity be added to light, beyond the previously mentioned characteristic,
namely that it annuls the concept and remains absolute being.
If we made
such additions, we would, as Kant has been criticized for doing, run up
against something unexplained and perhaps inexplicable.
As support for
this contention, notice that we have understood it as inconceivable purely
in its form and nothing more.
We have no right to assert anything before
we have seen into it.
So if we posit some other hidden quality,
we have either invented it, or better,
since pure invention from nothing is completely impossible,
we have manufactured it by trying to supply a principle for some facticity.
This happened with Kant when he first factically discovered
the distinction between the sensible and supersensible worlds and
then added to his absolute the additional inexplicable quality of linking the
two worlds, a move which pushed us back from genetic manifestness into
merely factical manifestness, completely contravening the inner spirit of
the science of knowing.
Therefore, it is important to note that whatever yet
to be determined characteristics the reality appearing in our knowing may
carry in itself, besides the common basic property of inconceivability, such
characteristics by no means require any new absolute grounding principle
besides the one principle of pure light, since this would multiply the num-
ber of absolutes.
Rather, the multiplicity and change of these various traits is to be
deduced purely from the interaction of the light with itself,
in its multifarious relations to concepts, and to inconceivable being.

I invite you to the following reflections
so that I can offer a hint about this last point,
raise what has been said today to a higher level,
give the new listeners a unified perspective
for viewing everything that they can learn here,
and give the returning listeners the same perspective
from which they can again gather and reproduce
everything they have heard up to now.

The focus of everything is pure light.
To truly come to this requires that
the concept be posited and annulled
and that an intrinsically inconceivable being be posited.
If it is granted that the light should exist,
then in this judgment everything else mentioned is possible as well.
We have now seen this; it is true; it remains true forever;
and it expresses the basic principle of all knowing.
We can so designate it for ourselves.

Now, however, I want to completely ignore the content of
this insight and reflect on its form, on our actual situation of insight.
I also think that we, those of us present here who have actually seen it,
are the ones who had the insight.
As I remember it and as I think we all do, the
process was that we freely constructed the concepts and premisses with
which we began, that we held them up to each other freely, and that in
holding them together we were gripped by the conviction that they
belonged together absolutely and formed an indivisible oneness.
Thus we created at least the conditions for the self-manifesting insight,
and so we likewise appeared to ourselves unconditionally.

But let us not go to work with too much haste;
rather let us consider things a little more deeply.
Did we create what we created because we wished to do so, and
therefore as the result of some earlier knowledge,
which we would have created because we wished to create it
as the result of an even earlier knowledge, and so on to infinity,
so that we might never arrive at a first creation?
Somewhere, if the concept is created it must absolutely and thoroughly
create itself, without anything antecedent and
without any necessity of a “we”;
because this “we,” as has been shown,
always and everywhere requires some previous knowing and
cannot achieve immediate knowing.
Therefore we cannot create the conditions,
they must emerge spontaneously.
Reason must create itself,
independent from any volition or freedom, or self.
But this proposition, disclosed through analysis,
[contradicts] the first which is given reflexively,
and so immediately.
Which one is true, and on which should we rely?
Before trying to answer, let us return again to the matter of the insight
which has become controversial in regard to its principle,
in order to grasp its meaning and true worth clearly.
We realized that if light is to be,
then the concept must be posited and negated.
Therefore, the light itself is not immediately present in this insight,
and the insight does not dissolve into light and coincide with it;
instead it is only an insight in relation to the light,
an insight which objectifies it,
grasping it by its inner quality only.
Thus, whatever the principle and true bearer of this insight might be,
whether we ourselves, as it seems to be,
or pure self-creating reason, as it also seems to be,
the light is not immediately present in this bearer,
instead it is present merely mediately in
a representative and likeness of itself.
First of all, that this light occurs merely mediately applies not
just in the science of knowing but in any possible consciousness
that has to posit a concept so it can annul it;
and the science of knowing rests on a completely different point
than the one on which many may have assumed it to rest after the last lecture,
because undoubtedly knowing was understood too simply.

And now to answer the question:
both are clear, therefore both are equally true;
and so, as was said at the start on another occasion,
manifestness rests neither in one nor in the other,
but entirely between both.
We arrive here, and this is the first important and significant result,
at the principle of division,
not as before a division between two terms,
which in that case are to be
intrinsically distinct like A and the point,
but instead a division of something
which remains always inwardly self-same
through all division.
In a word, we have to do with
constructing and creating
the very same primal concepts
which appear one time as immanent,
in the unconditionally evident final being, the I;
and appear the other time as emanent,
in reason, absolute and in itself,
which nevertheless is completely objectified.
Thus, it is a division pure and in itself,
without any result or alteration in the object.

Further, manifestness oscillates
between these two perspectives:
if it is to be really constructed,
then it must be constructed in that way.
Thus, it must be constructed as oscillating
from a to b and again from b to a and
as completely creating both;
thus as oscillating between the twofold oscillation,
which was the first point,
and which gives rise to a three or fivefold synthesis.

What is the common element in all these determinations?
The very same representative of the light,
seen in its familiar inner quality.
Here it stays.
Therefore everything is the same one common consciousness of light.
This consciousness, which is held in common and
therefore cannot be really constructed but
instead can only be thought by means of the science of knowing,
can be regarded or represented in the three or fivefold modifications only
from another standpoint, which this science alone oversees.
